\documentclass[11pt,a4paper]{article}
\usepackage[margin=1in]{geometry}
\usepackage[T1]{fontenc}
\usepackage[utf8]{inputenc}
\usepackage{lmodern}
\usepackage{microtype}
\usepackage{hyperref}
\usepackage{enumitem}
\setlist[itemize]{leftmargin=1.6em}
\setlist[enumerate]{leftmargin=1.8em}

\title{Why OpenMP Often Matches (Not Crushes) UPPAAL:\\
Chain Imbalance in Symbolic Reachability}
\author{PTCTCL Engineering Note}
\date{\today}

\begin{document}
\maketitle

\section{Executive Summary}
In this codebase, the OpenMP backend does \emph{symbolic} state-space exploration, not embarrassingly parallel per-state counting. Work is organized by discrete-state buckets (``chains''), where each chain holds an s antichain of zones and requires order-sensitive merge/subsumption logic.

That chain logic introduces:
\begin{itemize}
  \item unavoidable serial sections per chain,
  \item load imbalance across chains (a few heavy chains dominate level time), and
  \item synchronization overhead on shared structures.
\end{itemize}

So even with many cores, speedup is typically modest on small/medium models and only becomes visible on large models.

\section{What a ``Chain'' Is Here}
A chain corresponds to one discrete-state key $dk$ (locations + integer valuation) in the passed structure.
For a fixed $dk$, the checker stores multiple symbolic zones (DBMs) linked in one bucket/chain.

At each BFS level, new zones targeting the same $dk$ are compared against existing zones for:
\begin{enumerate}
  \item inclusion subsumption,
  \item reverse subsumption (killing old zones),
  \item convex union attempts, and
  \item insertion order constraints.
\end{enumerate}

This is the expensive part. Some $dk$ values stay tiny; a few become very large and dominate runtime.

\section{Why This Creates Serial Work}
Even with multiple threads, the per-chain update cannot be fully parallelized in a straightforward way because:
\begin{itemize}
  \item merge order matters for convex-union behavior and reproducibility,
  \item updates mutate chain membership (live/dead nodes),
  \item each candidate depends on previous survivors in that same chain.
\end{itemize}

So one thread may spend a long time on a heavy chain while others finish lighter chains and wait at the level barrier.

\section{A Small Example}
Assume one frontier level yields 8 chains with pending work sizes:
\[
[3, 4, 2, 5, 6, 4, 7, 200]
\]
With 8 worker threads, 7 threads quickly finish their small chains, but 1 thread handles the heavy chain (200). Overall level time is then close to:
\[
T_{level} \approx T_{heavy\ chain} + T_{barrier}
\]
not the average chain time. This is classic tail-latency imbalance.

A rough cost model for chain $i$:
\[
C_i \approx \alpha \cdot p_i \cdot l_i + \beta \cdot p_i^2
\]
where $p_i$ is pending zones and $l_i$ is current live antichain size.
Large $p_i$ or $l_i$ creates stragglers.

\section{Why More Threads Do Not Imply 32x Faster}
The intuition ``32 cores should process 32 states in the time 1 core processes 1'' fails because symbolic states are not uniform cheap units.

Performance is capped by:
\begin{itemize}
  \item serial fraction (Amdahl effect),
  \item lock/coherency overhead,
  \item memory bandwidth/cache pressure,
  \item chain-size skew (stragglers),
  \item barrier synchronization between levels.
\end{itemize}

So realistic speedups are often modest unless model size and structure are favorable.

\section{What Other Systems Do for Similar Problems}
This is a known issue in many domains (graph analytics, irregular sparse kernels, SAT/SMT search, symbolic state-space tools).
Common techniques:
\begin{enumerate}
  \item \textbf{Cost-aware scheduling}: estimate task cost and schedule heavy tasks first.
  \item \textbf{Work stealing}: idle workers steal remaining heavy tasks.
  \item \textbf{Chunked processing of heavy tasks}: split large chains into subchunks and requeue.
  \item \textbf{Adaptive policy switching}: choose incremental vs batched strategy based on online stats.
  \item \textbf{Hierarchical queues}: separate heavy and light queues to reduce tail effects.
  \item \textbf{NUMA-aware placement}: keep chain data close to the worker that processes it repeatedly.
\end{enumerate}

\section{How to Apply This Here (Practical Plan)}
\subsection*{Step 1: Instrumentation (first, no semantic change)}
Add per-level metrics:
\begin{itemize}
  \item number of touched chains,
  \item pending count per chain (max/median/p95),
  \item live antichain size per chain,
  \item settle time per chain,
  \item tail fraction: top 1\% chains time / total settle time.
\end{itemize}

If top 1\% dominates, load imbalance is confirmed as primary bottleneck.

\subsection*{Step 2: Cost-aware settle ordering}
Build task list $\{(dk, \hat C)\}$ with
\[
\hat C = p \cdot l + \gamma p^2
\]
and process descending by $\hat C$ with dynamic scheduling.

\subsection*{Step 3: Heavy-chain chunking}
For chains above threshold $\tau$, process only $K$ pending items per pass, then requeue.
This improves fairness and cuts barrier tail latency.

\subsection*{Step 4: Adaptive mode switch}
If pending-per-chain is small and contention low, prefer simple incremental path.
If skew/contension is high, switch to fused/settle-oriented path.

\section{Expected Outcome}
If chain skew is truly the bottleneck, these steps should:
\begin{itemize}
  \item reduce idle-at-barrier time,
  \item improve high-thread scaling on medium/large models,
  \item keep correctness (same reachability semantics) with only scheduling changes.
\end{itemize}

\section{Bottom Line}
Yes, the chain phenomenon is real and likely central in your current OpenMP behavior.
The right target is \emph{load-balance and tail reduction under chain-level serial constraints}, not just increasing thread count or tweaking chunk size globally.

\end{document}
